\section{Transferring the OPE estimates to homogeneous fields}
\label{sec:ope-estimates}
The raw prescription and its covariant homogeneous version need not agree
exactly on a product of fields. Their difference must therefore be estimated
before the homogeneous coefficients can be used in the expansion. We give
that transfer, retaining the complete correlations, including all unit and
disconnected terms. The quantitative estimate for the raw prescription
remains the construction input identified below.

Fix cluster fields $\boldsymbol v=(v_i)_{i\in I}$, with
$w=\sum_i\deg v_i$. For a sequence of positive numbers $r_n\to0$, set $t_n=-\log r_n\to\infty$.
Say that a sequence $f_n$ has upper rate $a$ if for every $\delta>0$,
\[
 \|f_n\|\leq e^{(a+\delta)t_n}=r_n^{-a-\delta}
                 \quad\hbox{for all sufficiently large }n.
\]
This notation permits zero values. Finite sums preserve a common upper
rate, products add rates, and fixed linear maps preserve rates.
For sums, first use half the desired slack and then absorb the finite
number of terms into $e^{\delta t_n/2}$. The same device absorbs a fixed
operator norm. Products follow by dividing the slack between the factors.
In particular any bounded sequence has upper rate zero.

If $\ell_n:E\to\mathbb C$ is a sequence of linear functionals on a
fixed finite-dimensional space, and $\ell_n(e)$ has upper rate $b$ for
every fixed $e$, then applying it to vectors $u_n$ of upper rate $a$
gives upper rate $a+b$. To prove this, expand $u_n$ in a fixed basis.
Each coordinate has rate at most $a$, each basis evaluation has rate at
most $b$, and the finite sum has rate at most $a+b$. The identical
finite-basis argument works for multilinear forms with several varying
inputs. We will use it for bounded rotated detector labels.

\begin{assumption}[Raw product comparison]\label{input:ope-raw-estimates}
For every fixed cluster label list, every $M\geq0$, $\eta>0$, every
positive sequence $r_n\to0$ and configurations with
\[
 |x_{i,n}-o_n|\leq Mr_n,\qquad
 |x_{i,n}-x_{j,n}|\geq\eta r_n\quad(i\neq j),
\]
the raw coefficients of \eqref{eq:ope-raw-coefficient} have upper rate
$w-|\alpha|$. For any fixed spectator label list $\boldsymbol z$,
any $c>0$, any positions satisfying
\[
 |y_{j,n}-o_n|\geq c,\qquad |y_{j,n}-y_{k,n}|\geq c\quad(j\neq k),
\]
and any sequence of phases or mixed-state choices $\chi_n$, the following
assertions hold.
For every $D$, the raw discrepancy
\begin{equation}\label{eq:ope-raw-discrepancy}
 S_{\chi_n}(\boldsymbol v,\boldsymbol z;\boldsymbol x_n,\boldsymbol y_n)
 -\ell_{\chi_n,D,o_n,\boldsymbol y_n,\boldsymbol z}
                 (H_D(\boldsymbol x_n,\boldsymbol v;o_n))
\end{equation}
has upper rate $w-D-1$. For every fixed $u\in E_j$, the separated
insertion $\ell_{\chi_n,j,o_n,\boldsymbol y_n,\boldsymbol z}(u)$ has
upper rate zero. These statements hold for every finite spectator list,
including empty lists, and for conjugated cluster fields as well.
The same underlying hierarchy and physical inclusions are used throughout.
\end{assumption}
In the source these are the fixed-exterior specialization of the raw
coefficient and full-discrepancy estimates; the stronger statements allow
several positive logarithmic scales. This input includes uniform control
along arbitrarily moving configurations and phases. A bound for each
fixed configuration with a threshold depending on that configuration
would be insufficient.

\subsection{The averaged readout error}
Fix a raw truncation degree $D$ and an output degree $m$, which may differ.
Abbreviate
$h_n=H_D(\boldsymbol x_n,\boldsymbol v;o_n)$. We claim that
\begin{equation}\label{eq:ope-average-bridge}
 C_m(\boldsymbol x_n,\boldsymbol v;o_n)-L_{mD}h_n
                 \quad\hbox{has upper rate }w-D-1.
\end{equation}
First consider only the rotation average. Let
\[
 Z_m(K;\boldsymbol x,\boldsymbol v;o)
       =R_m(K)H_m(K^{-1}\boldsymbol x,R(K^{-1})\boldsymbol v;K^{-1}o)
\]
be its integrand. On the common separated region, expanding the finite
tests gives the exact identity
\begin{align}
 &Z_m(K;\boldsymbol x,\boldsymbol v;o)
                  -Z_m(K;(o),(h);o)\nonumber\\
 &=\sum_{|\alpha|\leq m}\sum_\lambda b_{\alpha,\lambda}
   \bigl[S_{\mathrm{mix}}(\boldsymbol v,R(K)\boldsymbol w_\lambda;
                 \boldsymbol x,o+K\boldsymbol z_\lambda)
       -S_{\mathrm{mix}}(h,R(K)\boldsymbol w_\lambda;
                 o,o+K\boldsymbol z_\lambda)\bigr]
                         R_m(K)\iota_{|\alpha|,m}u_\alpha.
 \label{eq:ope-rotated-error}
\end{align}
Here each sum over $\lambda$ is restricted to the finite support of
$b_\alpha$. To verify the identity, write each raw coefficient as its
exterior correlation at relative origin. Rotate the entire correlation
by $K$ and translate by $o$; by the covariance and translation laws
\eqref{eq:ope-mixed-covariance} neither operation changes its value.
This is the only place where covariance of the correlations enters,
and it cannot be omitted: the assumptions on the actions $R_n$ and $J_n$
alone do not imply \eqref{eq:ope-rotated-error}. The cluster positions and labels become
$\boldsymbol x,\boldsymbol v$, and the detector positions and labels
become exactly those displayed. Subtract the same evaluation with the
single field $h$ at $o$, and use linearity in that slot.

Take $h=h_n$ and any sequence $K_n\in O(4)$. All the rotated detector
positions have a common positive separation from $o_n$ and from one
another: these are the fixed finite nonzero detector distances, preserved
by rotations. All rotated detector labels and output vectors are bounded,
by continuity of the representations on the compact group. Apply
Input~\ref{input:ope-raw-estimates} to each fixed basis choice of detector
labels in \eqref{eq:ope-rotated-error}, then use the multilinear
finite-basis argument above to allow their bounded variation with $K_n$.
Every summand and hence their finite sum has upper rate $w-D-1$.

This statement for every rotation sequence implies a uniform estimate in
$K$ at each sufficiently large $n$. Indeed, for a fixed slack, if there
were a violating rotation at infinitely many indices, choose one at each
such index and any rotation at the remaining indices. The resulting
sequence would contradict the rate estimate just proved. Thus the norm
of the integrand difference is eventually bounded uniformly by
$e^{(w-D-1+\delta)t_n}$. Both integrands are continuous on the compact
group once the shrinking cluster lies in $\mathcal U_m$. Integrating
against probability Haar measure gives the same bound for the difference
of their averages.

Apply the same argument to the conjugated part. Its single comparison
field is $J_Dh_n$, rather than the raw truncation recomputed from the
conjugated cluster. This distinction causes no extra assumption.
By the reality law in \eqref{eq:ope-mixed-covariance}, conjugating a
complete correlation conjugates all its labels, and the
raw discrepancy is linear in each spectator slot. Hence the discrepancy
with cluster $J\boldsymbol v$, comparison field $J_Dh_n$, and arbitrary
spectator labels is the complex conjugate of the original discrepancy
with conjugated spectator labels. That original discrepancy is covered
by the input and the bounded-label argument. It has the same norm and
upper rate. Output conjugation is a fixed real-linear map on a
finite-dimensional space and preserves the rate. Adding the two averages
and dividing by two proves \eqref{eq:ope-average-bridge}; translation
invariance \eqref{eq:ope-mixed-covariance} identifies its centered singleton readout with $L_{mD}h_n$.
Applying the fixed projection $K_m$ yields
\begin{equation}\label{eq:ope-fresh-bridge}
 F_m(\boldsymbol x_n,\boldsymbol v;o_n)-P_{mD}h_n
                    \quad\hbox{has upper rate }w-D-1.
\end{equation}

\subsection{Coefficient and remainder bounds}
Set $D=m$ in \eqref{eq:ope-fresh-bridge}. Applying $P_{mm}$ to the raw
shadow $H_m$ kills every term of degree less than $m$, by inclusion
compatibility and \eqref{eq:ope-higher-kills}. The remaining finitely many
terms consist of degree-$m$ raw coefficients, of upper rate $w-m$,
times fixed vectors. Their sum has that rate. The difference in
\eqref{eq:ope-fresh-bridge} has the smaller rate $w-m-1$, so
$F_m$ has upper rate $w-m$. Taking its fixed channel coordinates gives
the same rate for each coefficient in
\eqref{eq:ope-final-coefficients}.

For the remainder, retain the raw truncation $h_n=H_D$ and abbreviate
the spectator functional at degree $j$ by $\ell_{n,j}$. The exact
reconstruction identity in \eqref{eq:ope-projection-recursion} and
compatibility of spectator insertions give
\[
 \ell_{n,D}(h_n)=\sum_{j\leq D}\ell_{n,j}(P_{jD}h_n).
\]
Consequently the final homogeneous discrepancy is exactly the raw
discrepancy \eqref{eq:ope-raw-discrepancy} plus
\begin{equation}\label{eq:ope-final-correction}
 \sum_{j\leq D}\ell_{n,j}
       \bigl(P_{jD}h_n-F_j(\boldsymbol x_n,\boldsymbol v;o_n)\bigr).
\end{equation}
Each vector difference has upper rate $w-D-1$ by
\eqref{eq:ope-fresh-bridge}. The insertion functional has rate zero on
every fixed vector by the input. The finite-dimensional evaluation
argument therefore gives rate $w-D-1$ for every term of
\eqref{eq:ope-final-correction}, and for their finite sum. Adding the
raw discrepancy proves the same rate for the final discrepancy, along
any prescribed configuration and phase sequence.

\begin{theorem}\label{thm:ope-quantitative}
Under Input~\ref{input:ope-raw-estimates} and the standing assumptions of
Section~\ref{sec:ope-homogeneous}, including the covariance, translation and
reality laws \eqref{eq:ope-mixed-covariance} of the mixed correlations,
the coefficients of
\eqref{eq:ope-final-coefficients} satisfy the public separated OPE bounds.
Specifically, for every $\epsilon>0$, fixed channel $(m,\alpha)$ and
fixed $M\geq0$, $\eta>0$, at all sufficiently small $r>0$,
\[
 |c_{\boldsymbol v}^{m,\alpha}(\boldsymbol x;o)|
                      \leq r^{m-w-\epsilon}
\]
uniformly over $|x_i-o|\leq Mr$ and $|x_i-x_j|\geq\eta r$.
For each finite $D$, fixed spectator labels and $c>0$, one also has
\begin{align*}
 \bigg|S_\chi(\boldsymbol v,\boldsymbol z;\boldsymbol x,\boldsymbol y)
  -\sum_{j\leq D}\sum_\alpha
   c_{\boldsymbol v}^{j,\alpha}(\boldsymbol x;o)
   S_\chi(e_{j,\alpha},\boldsymbol z;o,\boldsymbol y)\bigg|
               \leq r^{D+1-w-\epsilon},
\end{align*}
uniformly over the same clusters, all spectators with center and mutual
clearance at least $c$, and all phases and the mixed state.
\end{theorem}
\begin{proof}
The preceding arguments prove these inequalities eventually along every
admissible sequence, by taking slack $\epsilon$ in the rate definition.
If a uniform threshold for the coefficient failed, choose for each $n$
a violating radius $0<r_n<1/n$ and an admissible cluster at that radius.
The resulting sequence contradicts its eventual rate bound. If the
uniform remainder threshold failed, choose in addition a violating
spectator configuration and a violating phase at each radius. The rate
bound was proved for arbitrary such choices, so gives the same
contradiction. A finite initial segment where a cluster might meet a
spectator is immaterial: for sufficiently small $r$, $Mr<c/2$ makes
the combined configuration separated, while the stated lower bounds
separate all points within either group. The center and spectators are
already separated by $c$.
\end{proof}
The coefficients used in this theorem were chosen before $D$, the phase
and the spectators. The small-radius threshold may depend on the cluster
and spectator labels, on $D$, $M$, $\eta$, $c$, $\epsilon$ and, for a
coefficient, on the channel; it does not depend on the phase, on the
mixed-state choice or on the positions. The theorem transfers the quantitative raw comparison;
it does not produce that comparison from the Wilson measure.
